% !PRINTED 139
\weberpage{139}

Accordingly we may confine ourselves to the consideration of \gesperrt{homogeneous symmetric functions}, from which all the others can be put together.

After this it is easy to give the general form of a symmetric function. We obtain it if, in one term of such a function,
\[
 \alpha_1^{\nu_1}\alpha_2^{\nu_2}\cdots\alpha_n^{\nu_n},
\]
as in the formation of determinants, we permute the lower indices in all possible ways and take the sum of all terms so formed. Such a function may be called an \gesperrt{elementary constituent} of a symmetric function. If we take several such elementary constituents, multiply them by arbitrary factors independent of the $\alpha$'s, and add them, then we obtain the most general symmetric function. The number of terms of one of these elementary constituents is $\Pi(n)$ when the exponents $\nu_1,\nu_2,\ldots,\nu_n$ are all different; but if one and the same exponent $\nu$ occurs several times, the permutations which give no different terms are to be omitted. For example, for three variables, if $\nu_1,\nu_2,\nu_3$ are different, an elementary constituent is
\[
\begin{aligned}
 &\alpha_1^{\nu_1}\alpha_2^{\nu_2}\alpha_3^{\nu_3}
+\alpha_1^{\nu_1}\alpha_3^{\nu_2}\alpha_2^{\nu_3}
+\alpha_2^{\nu_1}\alpha_1^{\nu_2}\alpha_3^{\nu_3}
+\alpha_2^{\nu_1}\alpha_3^{\nu_2}\alpha_1^{\nu_3}\\
 &\qquad+\alpha_3^{\nu_1}\alpha_1^{\nu_2}\alpha_2^{\nu_3}
+\alpha_3^{\nu_1}\alpha_2^{\nu_2}\alpha_1^{\nu_3},
\end{aligned}
\]
whereas, if $\nu_3=\nu_2$,
\[
 \alpha_1^{\nu_1}\alpha_2^{\nu_2}\alpha_3^{\nu_2}
+\alpha_2^{\nu_1}\alpha_1^{\nu_2}\alpha_3^{\nu_2}
+\alpha_3^{\nu_1}\alpha_1^{\nu_2}\alpha_2^{\nu_2}.
\]

The simplest example of a symmetric function is the sum of the variables
\[
 \alpha_1+\alpha_2+\alpha_3+\cdots+\alpha_n.
\]
Likewise the product $\alpha_1\alpha_2\cdots\alpha_n$ belongs here.

These two are the extreme cases of a series of symmetric functions which we call the \gesperrt{symmetric fundamental functions} and obtain as follows.

The product
\begin{equation}
 f(x)=(x-\alpha_1)(x-\alpha_2)\cdots(x-\alpha_n)
\tag{1}
\end{equation}
is, whatever $x$ may be, provided only that $x$ is independent of the $\alpha$'s, a symmetric function of $\alpha_1,\alpha_2,\ldots,\alpha_n$. If therefore we carry out the multiplication of the several factors and
%CONT>>
